\Article{15}
Es bleibt noch der Fall, dass der Punkt $\mu$ nicht auf der Fläche $t$ selbst liegt, sich aber in unmerklichem Abstand von ihr befindet. In diesem Fall kann unser Integral durchaus einen merklichen Wert haben, den wir jetzt genauer untersuchen wollen.

Die vom Punkt $\mu$ senkrecht zur Fläche $t$ gezogene Gerade und die vom selben Punkt ausgehende feste Gerade mögen die Kugeloberfläche in den Punkten $G$ beziehungsweise $H$ schneiden. Der Bogen sei $GH=k$, der Bogen zwischen $G$ und einem beliebigen Punkt der Kugeloberfläche sei $=v$. Schließlich sei $w$ der sphärische Winkel zwischen den Bögen $k$, $v$. Für das Element $\mathrm{d}\Pi$ können wir damit das Produkt $\sin v.\mathrm{d}v.\mathrm{d}w$ einsetzen. Schreiben wir der Kürze halber $r$ für $(\mu,\mathrm{d}t)$, so wird das Integral (II) zu
\[
=\int\int\pm(\cos k.\cos v+\sin k.\sin v.\cos w)\theta r.\sin v.\mathrm{d}v.\mathrm{d}w
\]
Diese Integration braucht nur über diejenigen Teile der Kugeloberfläche erstreckt zu werden, auf die die unmerklichen Abstände $r$ projiziert werden. Diese entsprechen dem unmerklichen Teil der Fläche $t$. Wenn wir ihn als \textit{eben} behandeln und den kleinsten Abstand, der zum Punkt $G$ oder zum Wert $v=0$ gehört, mit $\rho$ bezeichnen, ergibt sich $r=\frac{\rho}{\cos v}$, also unabhängig von $w$. Nach der Integration bezüglich der Variablen $w$, etwa von $w=0$ bis $w=360^0$, ergibt sich für das Integral
\[
=\pm\int2\pi\theta r.\cos k.\cos v.\sin v.\mathrm{d}v=\pm\int\frac{2\pi\cos k.\rho\rho\theta r.\mathrm{d}r}{r^3}
\]
Diese Integration ist von $r=\rho$ bis zu einem beliebigen, noch so kleinen merklichen Wert zu erstrecken. Setzen wir also allgemein
\[
2rr\int\frac{\theta r.\mathrm{d}r}{r^3}=-\theta'r
\]
\SourceLeaf{55}und wählen die Integrationskonstante so, dass $\int\frac{\theta r.\mathrm{d}r}{r^3}=0$ für einen beliebigen merklichen Wert in dem Bereich gilt, in dem wir Versuche anstellen können, so ergibt sich für das Integral (II) unter Vernachlässigung unmerklicher Größen
\[
=\pm\pi\cos k.\theta'\rho
\]

Erscheint es zweifelhaft, ob wir den Teil der Fläche $t$, der in unmerklichem Abstand vom Punkt $\mu$ liegt, als eben behandeln dürfen, so betrachten wir stattdessen eine Kugelfläche. Dabei sei $R$ der Abstand des Kugelmittelpunkts vom Punkt $\mu$, positiv oder negativ gerechnet, je nachdem, ob der Mittelpunkt in Richtung auf $G$ oder in der entgegengesetzten Richtung liegt. Dann gilt
\[
\begin{gathered}
\cos v=\frac{\rho}{r}(1-\frac{\rho}{2R})+\frac{r}{2R}\\
\sin v.\mathrm{d}v=[\frac{\rho}{rr}(1-\frac{\rho}{2R})-\frac{1}{2R}]\mathrm{d}r
\end{gathered}
\]
Daraus folgt leicht, dass sich das Integral in diesem Fall um keine merkliche Größe vom zuvor gefundenen Wert $\pm\pi\cos k.\theta'\rho$ unterscheidet, sofern $R$ eine merkliche Größe ist. Wie auch immer die Fläche $t$ im fraglichen Teil gekrümmt ist, lassen sich stets, sofern die Krümmungsradien merklich sind, zwei Kugelflächen angeben, die die Fläche $t$ an dem $\mu$ nächstgelegenen Punkt berühren, zwischen denen $t$ liegt und deren Radien von merklicher Größe sind. Offenbar liegt unser Integral dann zwischen den auf diese beiden Flächen bezogenen Integralen und wird somit ohne merklichen Fehler durch dieselbe Formel ausgedrückt. Eine Ausnahme tritt nur ein, wenn die Fläche $t$ in unmerklichem Abstand vom Punkt $\mu$ entweder eine Krümmung mit unmerklichem Radius, eine Kante oder eine Spitze aufweist.

\Article{16}
Gehen wir nun von der Integration (II) zum Integral (I) über, so wird dieses offenbar in zwei Fällen unmerklich. Der erste liegt vor, wenn jene für \textit{keinen} Punkt der Fläche $T$ einen merklichen Wert ergeben hat. Der zweite liegt vor, wenn die Gesamtheit der Elemente der Fläche $T$, für deren Punkte das Integral (II) merklich geworden ist, nur eine Fläche von unmerklicher Größe bildet. Bei genauer Betrachtung zeigt sich, dass das Integral (I) nur insoweit einen merklichen Wert annehmen kann, als die Fläche $T$ einen oder mehrere Teile von merklicher Größe enthält, die in unmerklichem Abstand von der Fläche $t$ liegen. Da solche Teile von der Parallelität zur Fläche $t$ nicht merklich abweichen können, wird sich für jeden ihrer Punkte $\cos k$ nicht merk\SourceLeaf{56}lich von $+1$ beziehungsweise $-1$ unterscheiden, je nachdem, ob die Fläche $T$ der Fläche $t$ ihre Außen- oder ihre Innenseite zuwendet. Bezeichnen wir also mit $\NearTau$, $\NearTau'$ diejenigen Teile der Fläche $T$, die in unmerklichem Abstand von der Fläche $t$ liegen, und zwar mit $\NearTau$ diejenigen, bei denen die Außenseite der einen Fläche der Innenseite der anderen zugewandt ist, und mit $\NearTau'$ diejenigen, bei denen gleichartige Seiten einander zugewandt sind, und bezeichnen wir schließlich mit $\rho$ den kleinsten Abstand eines beliebigen Elements $\mathrm{d}\NearTau$ oder $\mathrm{d}\NearTau'$ von der Fläche $t$, so wird unser Integral (I) unter Vernachlässigung unmerklicher Größen zu
\[
=-\int\pi\theta'\rho.\mathrm{d}\NearTau+\int\pi\theta'\rho.\mathrm{d}\NearTau'
\]
Offenbar macht es hier keinen Unterschied, ob die Teile $\NearTau$, $\NearTau'$ auf die Fläche $T$ oder auf $t$ bezogen werden.

Damit haben wir das in Artikel 6 gestellte Problem für diejenige Beschaffenheit der Funktion $\varphi$ vollständig gelöst, auf der unsere Hauptuntersuchung über die Gleichgewichtsform von Flüssigkeiten beruht. Wir erhalten
\[
\int\int\mathrm{d}s.\mathrm{d}S.\varphi(\mathrm{d}s,\mathrm{d}S)=4\pi\sigma\psi0-\pi\tau\theta0+\pi\tau'\theta0-\pi\int\mathrm{d}\NearTau.\theta'\rho+\pi\int\mathrm{d}\NearTau'.\theta'\rho
\]

\Article{17}
Wie die Funktion $\theta'$ entsteht, lässt sich durch die Gleichung
\[
\frac{\theta'r}{rr}=\int\frac{2\theta x.\mathrm{d}x}{x^3}
\]
ausdrücken, wobei das Integral von $x=r$ bis zu einem beliebigen konstanten merklichen Wert genommen wird, den wir hier mit $R$ bezeichnen. Offenbar ist dieses Integral kleiner als $\int\frac{2\theta r.\mathrm{d}x}{x^3}$ zwischen denselben Grenzen, dessen Wert $=\frac{\theta r}{rr}-\frac{\theta r}{RR}$ ist, und daher erst recht kleiner als $\frac{\theta r}{rr}$. Da aber für das unbestimmte Integral
\[
\int\frac{2\theta x.\mathrm{d}x}{x^3}=-\frac{\theta x}{xx}+\int\frac{\mathrm{d}\theta x}{xx}=-\frac{\theta x}{xx}-\int\frac{\psi x.\mathrm{d}x}{xx}
\]
gilt, erhalten wir
\[
\frac{\theta'r}{rr}=\frac{\theta r}{rr}-\frac{\theta R}{RR}-\int\frac{\psi x.\mathrm{d}x}{xx}
\]
wobei das Integral zwischen denselben Grenzen genommen wird. Dieses ist kleiner als das Integral $\int\frac{\psi r\mathrm{d}x}{xx}$ und daher auch kleiner als $\frac{\psi r}{r}$. Deshalb ist der Wert von $\frac{\theta'r}{rr}$ größer als
\[
\frac{\theta r}{rr}-\frac{\theta R}{RR}-\frac{\psi r}{r}
\]
\SourceLeaf{57}Somit liegt $\theta'r$ zwischen den Grenzen
\[
\theta r\quad\text{und}\quad\theta r-rr.\frac{\theta R}{RR}-r\psi r
\]
deren Unterschied bei unbegrenzt abnehmendem $r$ offenbar kleiner als jede angebbare Größe werden kann, da wir selbst $\psi0$ als endlich voraussetzen. Daraus schließen wir, dass $\theta'0=\theta0$ zu setzen ist. In der Formel, zu der wir im vorigen Artikel gelangt sind, lässt sich daher offenbar der Term $-\pi\tau\theta0$ als im Term $-\pi\int\mathrm{d}\NearTau.\theta'\rho$ enthalten ansehen und der Term $\pi\tau'\theta0$ als im Term $\pi\int\mathrm{d}\NearTau'.\theta'\rho$ enthalten, wenn wir den Unterschied zwischen einem unmerklichen Abstand und dem Abstand null aufheben und die Teile $\tau$, $\tau'$ jeweils zu den Teilen $\NearTau$, $\NearTau'$ hinzurechnen. Obwohl die Lösung dadurch im mathematischen Sinne eleganter würde, ist es für unseren Zweck jedoch besser, den Unterschied beizubehalten.

